\originalsection{有限导线与组合线源的实际积分}\label{tb:finite-b}
\begin{example}
\textbf{原9.1及9.11.1: 有限直导线.} 
电流$I$沿$x$轴从$x_1$到$x_2$流向$+x$, 求点$(x,y,0)$的磁场, $y\ne0$.
\end{example}
\begin{solution}
$\dd\vct\ell=\dd x'\uv x$, 位移$(x-x')\uv x+y\uv y$, 叉积只有$+z$分量$y\dd x'$. 因而
\[B_z=\frac{\mu_0I}{4\pi y}\left[\frac{x'-x}{\sqrt{(x'-x)^2+y^2}}\right]_{x_1}^{x_2}.\]
原例9.1在垂直平分线上, $x=0,x_1=-L/2,x_2=L/2$, 得
\[B_z=\frac{\mu_0I}{2\pi y}\frac{L/2}{\sqrt{(L/2)^2+y^2}}.\]
无限长极限为$\mu_0I/(2\pi y)$; 半无限线在端点的垂线上给一半. 写带符号的$y$可以自动确定两侧方向, 按大小写时另用右手法则. 用角度表示的端点公式与此等价, 但不同教材从导线或从垂线定义角度, 不可混用正弦与余弦.
\end{solution}

\begin{example}
\textbf{原9.11.2、9.11.4: 圆弧与发夹线.} 
(a)同心半径$a<b$的两段圆弧夹角$\theta_0$, 两端用径向线连接. 内弧电流顺时针、外弧逆时针, 求圆心场. (b)半径$r$的右半圆将两条向左伸到无穷远的水平导线连接, 电流先沿下线向右, 经弧到上线再向左, 求弧心场.
\end{example}
\begin{solution}
圆弧上$\dd\ell$与向心位移垂直, $\dd B=\mu_0I\dd\theta/(4\pi R)$. 径向线的叉积为零. (a)于是$B_z=\mu_0I\theta_0(1/b-1/a)/(4\pi)$, 指入纸面. (b)半圆贡献向外$\mu_0I/(4r)$, 两条半无限线各贡献向外$\mu_0I/(4\pi r)$, 总量
\[\bb=\frac{\mu_0I}{4\pi r}(\pi+2)\uv z.\]
三段方向相同是逐元右手方向的结果; 不能由于导线一进一出就以为两直线必相消.
\end{solution}

\begin{example}
\textbf{原9.11.3: 标题为矩形回路的U形无限线.} 
原图实际由三段构成: 从$(-a,+\infty)$向下到$(-a,b)$, 向右到$(a,b)$, 再向上到$(a,+\infty)$, $a,b>0$. 求原点场.
\end{example}
\begin{solution}
左右两竖段各给向外$\mu_0I[1-b/\sqrt{a^2+b^2}]/(4\pi a)$, 横段给向内$\mu_0Ia/[2\pi b\sqrt{a^2+b^2}]$. 相加得
\[B_z=\frac{\mu_0I}{2\pi ab}\left(b-\sqrt{a^2+b^2}\right)<0.\]
这是三段开放到无穷远的理想线模型, 无穷远的电流返回路径在本题中被忽略. 标题不能替代原图拓扑. $a\to0$时反向重合竖线相消、横段消失, 结果趋于零.
\end{solution}

\begin{example}
\textbf{原9.11.8: 直角折线.} 
电流从正$y$轴无穷远流向原点, 再沿正$x$轴流向无穷远. 求第一象限$(x,y)$的场.
\end{example}
\begin{solution}
竖段$\dd\vct\ell=-\uv y\dd s$, 横段$\dd\vct\ell=\uv x\dd s$, 两个叉积都向$+z$. 半无限公式给
\[B_z=\frac{\mu_0I}{4\pi}\left[\frac1x+\frac1y+\frac{y}{x\sqrt{x^2+y^2}}+\frac{x}{y\sqrt{x^2+y^2}}\right].\]
例如竖段来自$\int_0^\infty x\dd y'/[x^2+(y-y')^2]^{3/2}=(1+y/\sqrt{x^2+y^2})/x$. 两条线都对观测点有非零场, 不能将其中一条当作“径向段”删去.
\end{solution}

\originalsection{分布电流、运动电荷和转动圆盘}\label{tb:distributed-current}
\begin{example}
\textbf{原9.11.5--9.11.7: 三种非点线电流.} 
(a)两无限导线沿$-x$, 在$yz$面位置$(y,z)=(\pm a,0)$, 电流都为$I$. 求$z$轴场及最大场位置. (b)半径$R$圆柱内$J_z=\alpha s$. 求内外场. (c)宽$w$无限薄带内电流$I$沿$+x$均匀分布, 求同平面中带外、离近边$s$的场.
\end{example}
\begin{solution}
(a)各线环绕场相加, $z$分量相消, $B_y=\mu_0Iz/[\pi(a^2+z^2)]$. 导数为$\mu_0I(a^2-z^2)/[\pi(a^2+z^2)^2]$. 正半轴最大点$z=a$, $B=\mu_0I/(2\pi a)$; 原点场为零. 整个平面可写
\[B_y=\frac{\mu_0Iz}{2\pi}\sum_{\eta=\pm a}\frac1{(y-\eta)^2+z^2},\quad
B_z=-\frac{\mu_0I}{2\pi}\sum_{\eta=\pm a}\frac{y-\eta}{(y-\eta)^2+z^2}.\]
围单根线的场线顺时针, 远处两根合成总电流$2I$的环绕场. 本式给出了画全平面方向图的具体场, 原点零场不能替代整张场图.

\begin{center}\begin{tikzpicture}[>=Stealth,scale=.92]
\draw[->,gray] (-2.6,0)--(2.6,0) node[right] {$y/a$};\draw[->,gray] (0,-2.5)--(0,2.5) node[above] {$z/a$};
\draw[structurecolor] plot coordinates {(0.9762,-0.1240) (1.0046,-0.1244) (1.0167,-0.1226) (1.0259,-0.1204) (1.0336,-0.1180) (1.0404,-0.1155) (1.0464,-0.1128) (1.0519,-0.1100) (1.0569,-0.1072) (1.0616,-0.1043) (1.0659,-0.1013) (1.0700,-0.0983) (1.0737,-0.0952) (1.0773,-0.0921) (1.0806,-0.0889) (1.0838,-0.0858) (1.0867,-0.0825) (1.0895,-0.0793) (1.0921,-0.0760) (1.0946,-0.0727) (1.0969,-0.0694) (1.0990,-0.0660) (1.1011,-0.0627) (1.1030,-0.0593) (1.1048,-0.0559) (1.1064,-0.0525) (1.1080,-0.0490) (1.1094,-0.0456) (1.1107,-0.0421) (1.1119,-0.0386) (1.1130,-0.0352) (1.1139,-0.0317) (1.1148,-0.0282) (1.1156,-0.0247) (1.1162,-0.0212) (1.1168,-0.0176) (1.1172,-0.0141) (1.1176,-0.0106) (1.1178,-0.0071) (1.1180,-0.0035) (1.1180,0.0000) (1.1180,0.0035) (1.1178,0.0071) (1.1176,0.0106) (1.1172,0.0141) (1.1168,0.0176) (1.1162,0.0212) (1.1156,0.0247) (1.1148,0.0282) (1.1139,0.0317) (1.1130,0.0352) (1.1119,0.0386) (1.1107,0.0421) (1.1094,0.0456) (1.1080,0.0490) (1.1064,0.0525) (1.1048,0.0559) (1.1030,0.0593) (1.1011,0.0627) (1.0990,0.0660) (1.0969,0.0694) (1.0946,0.0727) (1.0921,0.0760) (1.0895,0.0793) (1.0867,0.0825) (1.0838,0.0858) (1.0806,0.0889) (1.0773,0.0921) (1.0737,0.0952) (1.0700,0.0983) (1.0659,0.1013) (1.0616,0.1043) (1.0569,0.1072) (1.0519,0.1100) (1.0464,0.1128) (1.0404,0.1155) (1.0336,0.1180) (1.0259,0.1204) (1.0167,0.1226) (1.0046,0.1244) (0.9762,0.1240) (0.9762,0.1240) (0.9492,0.1175) (0.9387,0.1132) (0.9310,0.1093) (0.9247,0.1056) (0.9194,0.1020) (0.9147,0.0986) (0.9105,0.0952) (0.9068,0.0920) (0.9034,0.0887) (0.9002,0.0856) (0.8974,0.0824) (0.8947,0.0793) (0.8923,0.0763) (0.8900,0.0733) (0.8879,0.0703) (0.8859,0.0673) (0.8840,0.0643) (0.8823,0.0614) (0.8807,0.0585) (0.8792,0.0556) (0.8778,0.0527) (0.8765,0.0499) (0.8753,0.0470) (0.8742,0.0442) (0.8732,0.0414) (0.8722,0.0386) (0.8713,0.0358) (0.8705,0.0330) (0.8698,0.0302) (0.8691,0.0275) (0.8685,0.0247) (0.8680,0.0219) (0.8675,0.0192) (0.8671,0.0164) (0.8668,0.0137) (0.8665,0.0109) (0.8663,0.0082) (0.8661,0.0055) (0.8661,0.0027) (0.8660,0.0000) (0.8661,-0.0027) (0.8661,-0.0055) (0.8663,-0.0082) (0.8665,-0.0109) (0.8668,-0.0137) (0.8671,-0.0164) (0.8675,-0.0192) (0.8680,-0.0219) (0.8685,-0.0247) (0.8691,-0.0275) (0.8698,-0.0302) (0.8705,-0.0330) (0.8713,-0.0358) (0.8722,-0.0386) (0.8732,-0.0414) (0.8742,-0.0442) (0.8753,-0.0470) (0.8765,-0.0499) (0.8778,-0.0527) (0.8792,-0.0556) (0.8807,-0.0585) (0.8823,-0.0614) (0.8840,-0.0643) (0.8859,-0.0673) (0.8879,-0.0703) (0.8900,-0.0733) (0.8923,-0.0763) (0.8947,-0.0793) (0.8974,-0.0824) (0.9002,-0.0856) (0.9034,-0.0887) (0.9068,-0.0920) (0.9105,-0.0952) (0.9147,-0.0986) (0.9194,-0.1020) (0.9247,-0.1056) (0.9310,-0.1093) (0.9387,-0.1132) (0.9492,-0.1175) (0.9762,-0.1240)}--cycle;
\draw[structurecolor] plot coordinates {(-0.9762,0.1240) (-1.0046,0.1244) (-1.0167,0.1226) (-1.0259,0.1204) (-1.0336,0.1180) (-1.0404,0.1155) (-1.0464,0.1128) (-1.0519,0.1100) (-1.0569,0.1072) (-1.0616,0.1043) (-1.0659,0.1013) (-1.0700,0.0983) (-1.0737,0.0952) (-1.0773,0.0921) (-1.0806,0.0889) (-1.0838,0.0858) (-1.0867,0.0825) (-1.0895,0.0793) (-1.0921,0.0760) (-1.0946,0.0727) (-1.0969,0.0694) (-1.0990,0.0660) (-1.1011,0.0627) (-1.1030,0.0593) (-1.1048,0.0559) (-1.1064,0.0525) (-1.1080,0.0490) (-1.1094,0.0456) (-1.1107,0.0421) (-1.1119,0.0386) (-1.1130,0.0352) (-1.1139,0.0317) (-1.1148,0.0282) (-1.1156,0.0247) (-1.1162,0.0212) (-1.1168,0.0176) (-1.1172,0.0141) (-1.1176,0.0106) (-1.1178,0.0071) (-1.1180,0.0035) (-1.1180,0.0000) (-1.1180,-0.0035) (-1.1178,-0.0071) (-1.1176,-0.0106) (-1.1172,-0.0141) (-1.1168,-0.0176) (-1.1162,-0.0212) (-1.1156,-0.0247) (-1.1148,-0.0282) (-1.1139,-0.0317) (-1.1130,-0.0352) (-1.1119,-0.0386) (-1.1107,-0.0421) (-1.1094,-0.0456) (-1.1080,-0.0490) (-1.1064,-0.0525) (-1.1048,-0.0559) (-1.1030,-0.0593) (-1.1011,-0.0627) (-1.0990,-0.0660) (-1.0969,-0.0694) (-1.0946,-0.0727) (-1.0921,-0.0760) (-1.0895,-0.0793) (-1.0867,-0.0825) (-1.0838,-0.0858) (-1.0806,-0.0889) (-1.0773,-0.0921) (-1.0737,-0.0952) (-1.0700,-0.0983) (-1.0659,-0.1013) (-1.0616,-0.1043) (-1.0569,-0.1072) (-1.0519,-0.1100) (-1.0464,-0.1128) (-1.0404,-0.1155) (-1.0336,-0.1180) (-1.0259,-0.1204) (-1.0167,-0.1226) (-1.0046,-0.1244) (-0.9762,-0.1240) (-0.9762,-0.1240) (-0.9492,-0.1175) (-0.9387,-0.1132) (-0.9310,-0.1093) (-0.9247,-0.1056) (-0.9194,-0.1020) (-0.9147,-0.0986) (-0.9105,-0.0952) (-0.9068,-0.0920) (-0.9034,-0.0887) (-0.9002,-0.0856) (-0.8974,-0.0824) (-0.8947,-0.0793) (-0.8923,-0.0763) (-0.8900,-0.0733) (-0.8879,-0.0703) (-0.8859,-0.0673) (-0.8840,-0.0643) (-0.8823,-0.0614) (-0.8807,-0.0585) (-0.8792,-0.0556) (-0.8778,-0.0527) (-0.8765,-0.0499) (-0.8753,-0.0470) (-0.8742,-0.0442) (-0.8732,-0.0414) (-0.8722,-0.0386) (-0.8713,-0.0358) (-0.8705,-0.0330) (-0.8698,-0.0302) (-0.8691,-0.0275) (-0.8685,-0.0247) (-0.8680,-0.0219) (-0.8675,-0.0192) (-0.8671,-0.0164) (-0.8668,-0.0137) (-0.8665,-0.0109) (-0.8663,-0.0082) (-0.8661,-0.0055) (-0.8661,-0.0027) (-0.8660,0.0000) (-0.8661,0.0027) (-0.8661,0.0055) (-0.8663,0.0082) (-0.8665,0.0109) (-0.8668,0.0137) (-0.8671,0.0164) (-0.8675,0.0192) (-0.8680,0.0219) (-0.8685,0.0247) (-0.8691,0.0275) (-0.8698,0.0302) (-0.8705,0.0330) (-0.8713,0.0358) (-0.8722,0.0386) (-0.8732,0.0414) (-0.8742,0.0442) (-0.8753,0.0470) (-0.8765,0.0499) (-0.8778,0.0527) (-0.8792,0.0556) (-0.8807,0.0585) (-0.8823,0.0614) (-0.8840,0.0643) (-0.8859,0.0673) (-0.8879,0.0703) (-0.8900,0.0733) (-0.8923,0.0763) (-0.8947,0.0793) (-0.8974,0.0824) (-0.9002,0.0856) (-0.9034,0.0887) (-0.9068,0.0920) (-0.9105,0.0952) (-0.9147,0.0986) (-0.9194,0.1020) (-0.9247,0.1056) (-0.9310,0.1093) (-0.9387,0.1132) (-0.9492,0.1175) (-0.9762,0.1240)}--cycle;
\draw[structurecolor] plot coordinates {(0.8754,-0.2624) (0.9398,-0.2742) (0.9684,-0.2749) (0.9909,-0.2735) (1.0102,-0.2710) (1.0273,-0.2676) (1.0429,-0.2635) (1.0573,-0.2590) (1.0706,-0.2540) (1.0831,-0.2487) (1.0948,-0.2430) (1.1058,-0.2370) (1.1163,-0.2307) (1.1261,-0.2242) (1.1354,-0.2175) (1.1443,-0.2106) (1.1527,-0.2035) (1.1606,-0.1962) (1.1681,-0.1887) (1.1752,-0.1811) (1.1819,-0.1733) (1.1883,-0.1654) (1.1943,-0.1574) (1.1999,-0.1492) (1.2051,-0.1410) (1.2100,-0.1327) (1.2146,-0.1242) (1.2188,-0.1157) (1.2228,-0.1071) (1.2263,-0.0984) (1.2296,-0.0897) (1.2325,-0.0809) (1.2352,-0.0720) (1.2375,-0.0631) (1.2395,-0.0542) (1.2412,-0.0452) (1.2425,-0.0362) (1.2436,-0.0272) (1.2444,-0.0181) (1.2448,-0.0091) (1.2450,0.0000) (1.2448,0.0091) (1.2444,0.0181) (1.2436,0.0272) (1.2425,0.0362) (1.2412,0.0452) (1.2395,0.0542) (1.2375,0.0631) (1.2352,0.0720) (1.2325,0.0809) (1.2296,0.0897) (1.2263,0.0984) (1.2228,0.1071) (1.2188,0.1157) (1.2146,0.1242) (1.2100,0.1327) (1.2051,0.1410) (1.1999,0.1492) (1.1943,0.1574) (1.1883,0.1654) (1.1819,0.1733) (1.1752,0.1811) (1.1681,0.1887) (1.1606,0.1962) (1.1527,0.2035) (1.1443,0.2106) (1.1354,0.2175) (1.1261,0.2242) (1.1163,0.2307) (1.1058,0.2370) (1.0948,0.2430) (1.0831,0.2487) (1.0706,0.2540) (1.0573,0.2590) (1.0429,0.2635) (1.0273,0.2676) (1.0102,0.2710) (0.9909,0.2735) (0.9684,0.2749) (0.9398,0.2742) (0.8754,0.2624) (0.8754,0.2624) (0.8190,0.2390) (0.7981,0.2266) (0.7832,0.2162) (0.7713,0.2069) (0.7613,0.1983) (0.7527,0.1902) (0.7452,0.1825) (0.7385,0.1752) (0.7325,0.1682) (0.7270,0.1613) (0.7221,0.1547) (0.7175,0.1483) (0.7134,0.1420) (0.7095,0.1359) (0.7059,0.1299) (0.7027,0.1240) (0.6996,0.1182) (0.6968,0.1126) (0.6942,0.1070) (0.6917,0.1014) (0.6895,0.0960) (0.6874,0.0906) (0.6854,0.0853) (0.6837,0.0800) (0.6820,0.0748) (0.6805,0.0696) (0.6791,0.0645) (0.6778,0.0594) (0.6767,0.0543) (0.6756,0.0493) (0.6747,0.0443) (0.6739,0.0393) (0.6731,0.0343) (0.6725,0.0294) (0.6720,0.0245) (0.6716,0.0196) (0.6712,0.0147) (0.6710,0.0098) (0.6709,0.0049) (0.6708,0.0000) (0.6709,-0.0049) (0.6710,-0.0098) (0.6712,-0.0147) (0.6716,-0.0196) (0.6720,-0.0245) (0.6725,-0.0294) (0.6731,-0.0343) (0.6739,-0.0393) (0.6747,-0.0443) (0.6756,-0.0493) (0.6767,-0.0543) (0.6778,-0.0594) (0.6791,-0.0645) (0.6805,-0.0696) (0.6820,-0.0748) (0.6837,-0.0800) (0.6854,-0.0853) (0.6874,-0.0906) (0.6895,-0.0960) (0.6917,-0.1014) (0.6942,-0.1070) (0.6968,-0.1126) (0.6996,-0.1182) (0.7027,-0.1240) (0.7059,-0.1299) (0.7095,-0.1359) (0.7134,-0.1420) (0.7175,-0.1483) (0.7221,-0.1547) (0.7270,-0.1613) (0.7325,-0.1682) (0.7385,-0.1752) (0.7452,-0.1825) (0.7527,-0.1902) (0.7613,-0.1983) (0.7713,-0.2069) (0.7832,-0.2162) (0.7981,-0.2266) (0.8190,-0.2390) (0.8754,-0.2624)}--cycle;
\draw[structurecolor] plot coordinates {(-0.8754,0.2624) (-0.9398,0.2742) (-0.9684,0.2749) (-0.9909,0.2735) (-1.0102,0.2710) (-1.0273,0.2676) (-1.0429,0.2635) (-1.0573,0.2590) (-1.0706,0.2540) (-1.0831,0.2487) (-1.0948,0.2430) (-1.1058,0.2370) (-1.1163,0.2307) (-1.1261,0.2242) (-1.1354,0.2175) (-1.1443,0.2106) (-1.1527,0.2035) (-1.1606,0.1962) (-1.1681,0.1887) (-1.1752,0.1811) (-1.1819,0.1733) (-1.1883,0.1654) (-1.1943,0.1574) (-1.1999,0.1492) (-1.2051,0.1410) (-1.2100,0.1327) (-1.2146,0.1242) (-1.2188,0.1157) (-1.2228,0.1071) (-1.2263,0.0984) (-1.2296,0.0897) (-1.2325,0.0809) (-1.2352,0.0720) (-1.2375,0.0631) (-1.2395,0.0542) (-1.2412,0.0452) (-1.2425,0.0362) (-1.2436,0.0272) (-1.2444,0.0181) (-1.2448,0.0091) (-1.2450,0.0000) (-1.2448,-0.0091) (-1.2444,-0.0181) (-1.2436,-0.0272) (-1.2425,-0.0362) (-1.2412,-0.0452) (-1.2395,-0.0542) (-1.2375,-0.0631) (-1.2352,-0.0720) (-1.2325,-0.0809) (-1.2296,-0.0897) (-1.2263,-0.0984) (-1.2228,-0.1071) (-1.2188,-0.1157) (-1.2146,-0.1242) (-1.2100,-0.1327) (-1.2051,-0.1410) (-1.1999,-0.1492) (-1.1943,-0.1574) (-1.1883,-0.1654) (-1.1819,-0.1733) (-1.1752,-0.1811) (-1.1681,-0.1887) (-1.1606,-0.1962) (-1.1527,-0.2035) (-1.1443,-0.2106) (-1.1354,-0.2175) (-1.1261,-0.2242) (-1.1163,-0.2307) (-1.1058,-0.2370) (-1.0948,-0.2430) (-1.0831,-0.2487) (-1.0706,-0.2540) (-1.0573,-0.2590) (-1.0429,-0.2635) (-1.0273,-0.2676) (-1.0102,-0.2710) (-0.9909,-0.2735) (-0.9684,-0.2749) (-0.9398,-0.2742) (-0.8754,-0.2624) (-0.8754,-0.2624) (-0.8190,-0.2390) (-0.7981,-0.2266) (-0.7832,-0.2162) (-0.7713,-0.2069) (-0.7613,-0.1983) (-0.7527,-0.1902) (-0.7452,-0.1825) (-0.7385,-0.1752) (-0.7325,-0.1682) (-0.7270,-0.1613) (-0.7221,-0.1547) (-0.7175,-0.1483) (-0.7134,-0.1420) (-0.7095,-0.1359) (-0.7059,-0.1299) (-0.7027,-0.1240) (-0.6996,-0.1182) (-0.6968,-0.1126) (-0.6942,-0.1070) (-0.6917,-0.1014) (-0.6895,-0.0960) (-0.6874,-0.0906) (-0.6854,-0.0853) (-0.6837,-0.0800) (-0.6820,-0.0748) (-0.6805,-0.0696) (-0.6791,-0.0645) (-0.6778,-0.0594) (-0.6767,-0.0543) (-0.6756,-0.0493) (-0.6747,-0.0443) (-0.6739,-0.0393) (-0.6731,-0.0343) (-0.6725,-0.0294) (-0.6720,-0.0245) (-0.6716,-0.0196) (-0.6712,-0.0147) (-0.6710,-0.0098) (-0.6709,-0.0049) (-0.6708,0.0000) (-0.6709,0.0049) (-0.6710,0.0098) (-0.6712,0.0147) (-0.6716,0.0196) (-0.6720,0.0245) (-0.6725,0.0294) (-0.6731,0.0343) (-0.6739,0.0393) (-0.6747,0.0443) (-0.6756,0.0493) (-0.6767,0.0543) (-0.6778,0.0594) (-0.6791,0.0645) (-0.6805,0.0696) (-0.6820,0.0748) (-0.6837,0.0800) (-0.6854,0.0853) (-0.6874,0.0906) (-0.6895,0.0960) (-0.6917,0.1014) (-0.6942,0.1070) (-0.6968,0.1126) (-0.6996,0.1182) (-0.7027,0.1240) (-0.7059,0.1299) (-0.7095,0.1359) (-0.7134,0.1420) (-0.7175,0.1483) (-0.7221,0.1547) (-0.7270,0.1613) (-0.7325,0.1682) (-0.7385,0.1752) (-0.7452,0.1825) (-0.7527,0.1902) (-0.7613,0.1983) (-0.7713,0.2069) (-0.7832,0.2162) (-0.7981,0.2266) (-0.8190,0.2390) (-0.8754,0.2624)}--cycle;
\draw[structurecolor] plot coordinates {(0.6341,-0.3530) (0.7361,-0.3976) (0.7848,-0.4112) (0.8247,-0.4188) (0.8599,-0.4230) (0.8920,-0.4248) (0.9217,-0.4247) (0.9497,-0.4231) (0.9761,-0.4201) (1.0012,-0.4159) (1.0251,-0.4106) (1.0479,-0.4044) (1.0697,-0.3973) (1.0906,-0.3893) (1.1105,-0.3806) (1.1296,-0.3712) (1.1479,-0.3611) (1.1654,-0.3504) (1.1821,-0.3391) (1.1980,-0.3273) (1.2131,-0.3149) (1.2275,-0.3021) (1.2411,-0.2888) (1.2540,-0.2750) (1.2661,-0.2609) (1.2775,-0.2463) (1.2881,-0.2315) (1.2981,-0.2163) (1.3072,-0.2008) (1.3157,-0.1850) (1.3234,-0.1690) (1.3304,-0.1527) (1.3366,-0.1363) (1.3421,-0.1196) (1.3469,-0.1028) (1.3510,-0.0859) (1.3543,-0.0689) (1.3568,-0.0517) (1.3587,-0.0345) (1.3598,-0.0173) (1.3601,0.0000) (1.3598,0.0173) (1.3587,0.0345) (1.3568,0.0517) (1.3543,0.0689) (1.3510,0.0859) (1.3469,0.1028) (1.3421,0.1196) (1.3366,0.1363) (1.3304,0.1527) (1.3234,0.1690) (1.3157,0.1850) (1.3072,0.2008) (1.2981,0.2163) (1.2881,0.2315) (1.2775,0.2463) (1.2661,0.2609) (1.2540,0.2750) (1.2411,0.2888) (1.2275,0.3021) (1.2131,0.3149) (1.1980,0.3273) (1.1821,0.3391) (1.1654,0.3504) (1.1479,0.3611) (1.1296,0.3712) (1.1105,0.3806) (1.0906,0.3893) (1.0697,0.3973) (1.0479,0.4044) (1.0251,0.4106) (1.0012,0.4159) (0.9761,0.4201) (0.9497,0.4231) (0.9217,0.4247) (0.8920,0.4248) (0.8599,0.4230) (0.8247,0.4188) (0.7848,0.4112) (0.7361,0.3976) (0.6341,0.3530) (0.6341,0.3530) (0.5540,0.2993) (0.5266,0.2759) (0.5078,0.2579) (0.4932,0.2426) (0.4814,0.2293) (0.4714,0.2172) (0.4628,0.2062) (0.4554,0.1960) (0.4487,0.1864) (0.4428,0.1774) (0.4376,0.1689) (0.4328,0.1607) (0.4284,0.1530) (0.4245,0.1455) (0.4209,0.1383) (0.4176,0.1314) (0.4145,0.1247) (0.4118,0.1181) (0.4092,0.1118) (0.4068,0.1056) (0.4047,0.0996) (0.4027,0.0937) (0.4008,0.0879) (0.3991,0.0822) (0.3976,0.0767) (0.3962,0.0712) (0.3949,0.0658) (0.3937,0.0605) (0.3926,0.0552) (0.3917,0.0500) (0.3908,0.0449) (0.3901,0.0398) (0.3894,0.0347) (0.3888,0.0297) (0.3884,0.0247) (0.3880,0.0197) (0.3877,0.0148) (0.3875,0.0098) (0.3873,0.0049) (0.3873,0.0000) (0.3873,-0.0049) (0.3875,-0.0098) (0.3877,-0.0148) (0.3880,-0.0197) (0.3884,-0.0247) (0.3888,-0.0297) (0.3894,-0.0347) (0.3901,-0.0398) (0.3908,-0.0449) (0.3917,-0.0500) (0.3926,-0.0552) (0.3937,-0.0605) (0.3949,-0.0658) (0.3962,-0.0712) (0.3976,-0.0767) (0.3991,-0.0822) (0.4008,-0.0879) (0.4027,-0.0937) (0.4047,-0.0996) (0.4068,-0.1056) (0.4092,-0.1118) (0.4118,-0.1181) (0.4145,-0.1247) (0.4176,-0.1314) (0.4209,-0.1383) (0.4245,-0.1455) (0.4284,-0.1530) (0.4328,-0.1607) (0.4376,-0.1689) (0.4428,-0.1774) (0.4487,-0.1864) (0.4554,-0.1960) (0.4628,-0.2062) (0.4714,-0.2172) (0.4814,-0.2293) (0.4932,-0.2426) (0.5078,-0.2579) (0.5266,-0.2759) (0.5540,-0.2993) (0.6341,-0.3530)}--cycle;
\draw[structurecolor] plot coordinates {(-0.6341,0.3530) (-0.7361,0.3976) (-0.7848,0.4112) (-0.8247,0.4188) (-0.8599,0.4230) (-0.8920,0.4248) (-0.9217,0.4247) (-0.9497,0.4231) (-0.9761,0.4201) (-1.0012,0.4159) (-1.0251,0.4106) (-1.0479,0.4044) (-1.0697,0.3973) (-1.0906,0.3893) (-1.1105,0.3806) (-1.1296,0.3712) (-1.1479,0.3611) (-1.1654,0.3504) (-1.1821,0.3391) (-1.1980,0.3273) (-1.2131,0.3149) (-1.2275,0.3021) (-1.2411,0.2888) (-1.2540,0.2750) (-1.2661,0.2609) (-1.2775,0.2463) (-1.2881,0.2315) (-1.2981,0.2163) (-1.3072,0.2008) (-1.3157,0.1850) (-1.3234,0.1690) (-1.3304,0.1527) (-1.3366,0.1363) (-1.3421,0.1196) (-1.3469,0.1028) (-1.3510,0.0859) (-1.3543,0.0689) (-1.3568,0.0517) (-1.3587,0.0345) (-1.3598,0.0173) (-1.3601,0.0000) (-1.3598,-0.0173) (-1.3587,-0.0345) (-1.3568,-0.0517) (-1.3543,-0.0689) (-1.3510,-0.0859) (-1.3469,-0.1028) (-1.3421,-0.1196) (-1.3366,-0.1363) (-1.3304,-0.1527) (-1.3234,-0.1690) (-1.3157,-0.1850) (-1.3072,-0.2008) (-1.2981,-0.2163) (-1.2881,-0.2315) (-1.2775,-0.2463) (-1.2661,-0.2609) (-1.2540,-0.2750) (-1.2411,-0.2888) (-1.2275,-0.3021) (-1.2131,-0.3149) (-1.1980,-0.3273) (-1.1821,-0.3391) (-1.1654,-0.3504) (-1.1479,-0.3611) (-1.1296,-0.3712) (-1.1105,-0.3806) (-1.0906,-0.3893) (-1.0697,-0.3973) (-1.0479,-0.4044) (-1.0251,-0.4106) (-1.0012,-0.4159) (-0.9761,-0.4201) (-0.9497,-0.4231) (-0.9217,-0.4247) (-0.8920,-0.4248) (-0.8599,-0.4230) (-0.8247,-0.4188) (-0.7848,-0.4112) (-0.7361,-0.3976) (-0.6341,-0.3530) (-0.6341,-0.3530) (-0.5540,-0.2993) (-0.5266,-0.2759) (-0.5078,-0.2579) (-0.4932,-0.2426) (-0.4814,-0.2293) (-0.4714,-0.2172) (-0.4628,-0.2062) (-0.4554,-0.1960) (-0.4487,-0.1864) (-0.4428,-0.1774) (-0.4376,-0.1689) (-0.4328,-0.1607) (-0.4284,-0.1530) (-0.4245,-0.1455) (-0.4209,-0.1383) (-0.4176,-0.1314) (-0.4145,-0.1247) (-0.4118,-0.1181) (-0.4092,-0.1118) (-0.4068,-0.1056) (-0.4047,-0.0996) (-0.4027,-0.0937) (-0.4008,-0.0879) (-0.3991,-0.0822) (-0.3976,-0.0767) (-0.3962,-0.0712) (-0.3949,-0.0658) (-0.3937,-0.0605) (-0.3926,-0.0552) (-0.3917,-0.0500) (-0.3908,-0.0449) (-0.3901,-0.0398) (-0.3894,-0.0347) (-0.3888,-0.0297) (-0.3884,-0.0247) (-0.3880,-0.0197) (-0.3877,-0.0148) (-0.3875,-0.0098) (-0.3873,-0.0049) (-0.3873,0.0000) (-0.3873,0.0049) (-0.3875,0.0098) (-0.3877,0.0148) (-0.3880,0.0197) (-0.3884,0.0247) (-0.3888,0.0297) (-0.3894,0.0347) (-0.3901,0.0398) (-0.3908,0.0449) (-0.3917,0.0500) (-0.3926,0.0552) (-0.3937,0.0605) (-0.3949,0.0658) (-0.3962,0.0712) (-0.3976,0.0767) (-0.3991,0.0822) (-0.4008,0.0879) (-0.4027,0.0937) (-0.4047,0.0996) (-0.4068,0.1056) (-0.4092,0.1118) (-0.4118,0.1181) (-0.4145,0.1247) (-0.4176,0.1314) (-0.4209,0.1383) (-0.4245,0.1455) (-0.4284,0.1530) (-0.4328,0.1607) (-0.4376,0.1689) (-0.4428,0.1774) (-0.4487,0.1864) (-0.4554,0.1960) (-0.4628,0.2062) (-0.4714,0.2172) (-0.4814,0.2293) (-0.4932,0.2426) (-0.5078,0.2579) (-0.5266,0.2759) (-0.5540,0.2993) (-0.6341,0.3530)}--cycle;
\draw[structurecolor] plot coordinates {(1.4663,0.0000) (1.4643,0.0460) (1.4584,0.0918) (1.4485,0.1369) (1.4348,0.1813) (1.4172,0.2245) (1.3958,0.2663) (1.3708,0.3064) (1.3423,0.3446) (1.3103,0.3807) (1.2751,0.4143) (1.2367,0.4452) (1.1954,0.4733) (1.1514,0.4983) (1.1049,0.5199) (1.0562,0.5381) (1.0055,0.5528) (0.9532,0.5637) (0.8997,0.5710) (0.8454,0.5746) (0.7908,0.5746) (0.7364,0.5712) (0.6828,0.5649) (0.6307,0.5560) (0.5805,0.5451) (0.5329,0.5329) (0.4882,0.5199) (0.4467,0.5067) (0.4085,0.4937) (0.3735,0.4815) (0.3415,0.4700) (0.3123,0.4595) (0.2856,0.4500) (0.2610,0.4414) (0.2384,0.4337) (0.2175,0.4269) (0.1980,0.4208) (0.1798,0.4154) (0.1626,0.4107) (0.1464,0.4065) (0.1309,0.4028) (0.1161,0.3997) (0.1019,0.3969) (0.0882,0.3946) (0.0749,0.3926) (0.0619,0.3909) (0.0492,0.3896) (0.0367,0.3886) (0.0244,0.3879) (0.0122,0.3874) (0.0000,0.3873) (-0.0122,0.3874) (-0.0244,0.3879) (-0.0367,0.3886) (-0.0492,0.3896) (-0.0619,0.3909) (-0.0749,0.3926) (-0.0882,0.3946) (-0.1019,0.3969) (-0.1161,0.3997) (-0.1309,0.4028) (-0.1464,0.4065) (-0.1626,0.4107) (-0.1798,0.4154) (-0.1980,0.4208) (-0.2175,0.4269) (-0.2384,0.4337) (-0.2610,0.4414) (-0.2856,0.4500) (-0.3123,0.4595) (-0.3415,0.4700) (-0.3735,0.4815) (-0.4085,0.4937) (-0.4467,0.5067) (-0.4882,0.5199) (-0.5329,0.5329) (-0.5805,0.5451) (-0.6307,0.5560) (-0.6828,0.5649) (-0.7364,0.5712) (-0.7908,0.5746) (-0.8454,0.5746) (-0.8997,0.5710) (-0.9532,0.5637) (-1.0055,0.5528) (-1.0562,0.5381) (-1.1049,0.5199) (-1.1514,0.4983) (-1.1954,0.4733) (-1.2367,0.4452) (-1.2751,0.4143) (-1.3103,0.3807) (-1.3423,0.3446) (-1.3708,0.3064) (-1.3958,0.2663) (-1.4172,0.2245) (-1.4348,0.1813) (-1.4485,0.1369) (-1.4584,0.0918) (-1.4643,0.0460) (-1.4663,0.0000) (-1.4643,-0.0460) (-1.4584,-0.0918) (-1.4485,-0.1369) (-1.4348,-0.1813) (-1.4172,-0.2245) (-1.3958,-0.2663) (-1.3708,-0.3064) (-1.3423,-0.3446) (-1.3103,-0.3807) (-1.2751,-0.4143) (-1.2367,-0.4452) (-1.1954,-0.4733) (-1.1514,-0.4983) (-1.1049,-0.5199) (-1.0562,-0.5381) (-1.0055,-0.5528) (-0.9532,-0.5637) (-0.8997,-0.5710) (-0.8454,-0.5746) (-0.7908,-0.5746) (-0.7364,-0.5712) (-0.6828,-0.5649) (-0.6307,-0.5560) (-0.5805,-0.5451) (-0.5329,-0.5329) (-0.4882,-0.5199) (-0.4467,-0.5067) (-0.4085,-0.4937) (-0.3735,-0.4815) (-0.3415,-0.4700) (-0.3123,-0.4595) (-0.2856,-0.4500) (-0.2610,-0.4414) (-0.2384,-0.4337) (-0.2175,-0.4269) (-0.1980,-0.4208) (-0.1798,-0.4154) (-0.1626,-0.4107) (-0.1464,-0.4065) (-0.1309,-0.4028) (-0.1161,-0.3997) (-0.1019,-0.3969) (-0.0882,-0.3946) (-0.0749,-0.3926) (-0.0619,-0.3909) (-0.0492,-0.3896) (-0.0367,-0.3886) (-0.0244,-0.3879) (-0.0122,-0.3874) (-0.0000,-0.3873) (0.0122,-0.3874) (0.0244,-0.3879) (0.0367,-0.3886) (0.0492,-0.3896) (0.0619,-0.3909) (0.0749,-0.3926) (0.0882,-0.3946) (0.1019,-0.3969) (0.1161,-0.3997) (0.1309,-0.4028) (0.1464,-0.4065) (0.1626,-0.4107) (0.1798,-0.4154) (0.1980,-0.4208) (0.2175,-0.4269) (0.2384,-0.4337) (0.2610,-0.4414) (0.2856,-0.4500) (0.3123,-0.4595) (0.3415,-0.4700) (0.3735,-0.4815) (0.4085,-0.4937) (0.4467,-0.5067) (0.4882,-0.5199) (0.5329,-0.5329) (0.5805,-0.5451) (0.6307,-0.5560) (0.6828,-0.5649) (0.7364,-0.5712) (0.7908,-0.5746) (0.8454,-0.5746) (0.8997,-0.5710) (0.9532,-0.5637) (1.0055,-0.5528) (1.0562,-0.5381) (1.1049,-0.5199) (1.1514,-0.4983) (1.1954,-0.4733) (1.2367,-0.4452) (1.2751,-0.4143) (1.3103,-0.3807) (1.3423,-0.3446) (1.3708,-0.3064) (1.3958,-0.2663) (1.4172,-0.2245) (1.4348,-0.1813) (1.4485,-0.1369) (1.4584,-0.0918) (1.4643,-0.0460) (1.4663,-0.0000)};
\draw[structurecolor] plot coordinates {(1.6125,0.0000) (1.6107,0.0506) (1.6053,0.1010) (1.5964,0.1509) (1.5840,0.2001) (1.5682,0.2484) (1.5490,0.2955) (1.5267,0.3412) (1.5012,0.3854) (1.4727,0.4279) (1.4415,0.4684) (1.4076,0.5068) (1.3713,0.5429) (1.3327,0.5767) (1.2922,0.6081) (1.2500,0.6369) (1.2063,0.6632) (1.1613,0.6868) (1.1155,0.7079) (1.0689,0.7264) (1.0220,0.7425) (0.9749,0.7562) (0.9279,0.7677) (0.8814,0.7770) (0.8354,0.7845) (0.7903,0.7903) (0.7461,0.7945) (0.7030,0.7974) (0.6611,0.7991) (0.6205,0.7999) (0.5812,0.7999) (0.5432,0.7993) (0.5066,0.7982) (0.4712,0.7968) (0.4371,0.7951) (0.4042,0.7933) (0.3724,0.7913) (0.3416,0.7893) (0.3118,0.7874) (0.2828,0.7855) (0.2547,0.7837) (0.2272,0.7821) (0.2004,0.7806) (0.1742,0.7792) (0.1484,0.7780) (0.1231,0.7770) (0.0980,0.7761) (0.0733,0.7755) (0.0488,0.7750) (0.0243,0.7747) (0.0000,0.7746) (-0.0243,0.7747) (-0.0488,0.7750) (-0.0733,0.7755) (-0.0980,0.7761) (-0.1231,0.7770) (-0.1484,0.7780) (-0.1742,0.7792) (-0.2004,0.7806) (-0.2272,0.7821) (-0.2547,0.7837) (-0.2828,0.7855) (-0.3118,0.7874) (-0.3416,0.7893) (-0.3724,0.7913) (-0.4042,0.7933) (-0.4371,0.7951) (-0.4712,0.7968) (-0.5066,0.7982) (-0.5432,0.7993) (-0.5812,0.7999) (-0.6205,0.7999) (-0.6611,0.7991) (-0.7030,0.7974) (-0.7461,0.7945) (-0.7903,0.7903) (-0.8354,0.7845) (-0.8814,0.7770) (-0.9279,0.7677) (-0.9749,0.7562) (-1.0220,0.7425) (-1.0689,0.7264) (-1.1155,0.7079) (-1.1613,0.6868) (-1.2063,0.6632) (-1.2500,0.6369) (-1.2922,0.6081) (-1.3327,0.5767) (-1.3713,0.5429) (-1.4076,0.5068) (-1.4415,0.4684) (-1.4727,0.4279) (-1.5012,0.3854) (-1.5267,0.3412) (-1.5490,0.2955) (-1.5682,0.2484) (-1.5840,0.2001) (-1.5964,0.1509) (-1.6053,0.1010) (-1.6107,0.0506) (-1.6125,0.0000) (-1.6107,-0.0506) (-1.6053,-0.1010) (-1.5964,-0.1509) (-1.5840,-0.2001) (-1.5682,-0.2484) (-1.5490,-0.2955) (-1.5267,-0.3412) (-1.5012,-0.3854) (-1.4727,-0.4279) (-1.4415,-0.4684) (-1.4076,-0.5068) (-1.3713,-0.5429) (-1.3327,-0.5767) (-1.2922,-0.6081) (-1.2500,-0.6369) (-1.2063,-0.6632) (-1.1613,-0.6868) (-1.1155,-0.7079) (-1.0689,-0.7264) (-1.0220,-0.7425) (-0.9749,-0.7562) (-0.9279,-0.7677) (-0.8814,-0.7770) (-0.8354,-0.7845) (-0.7903,-0.7903) (-0.7461,-0.7945) (-0.7030,-0.7974) (-0.6611,-0.7991) (-0.6205,-0.7999) (-0.5812,-0.7999) (-0.5432,-0.7993) (-0.5066,-0.7982) (-0.4712,-0.7968) (-0.4371,-0.7951) (-0.4042,-0.7933) (-0.3724,-0.7913) (-0.3416,-0.7893) (-0.3118,-0.7874) (-0.2828,-0.7855) (-0.2547,-0.7837) (-0.2272,-0.7821) (-0.2004,-0.7806) (-0.1742,-0.7792) (-0.1484,-0.7780) (-0.1231,-0.7770) (-0.0980,-0.7761) (-0.0733,-0.7755) (-0.0488,-0.7750) (-0.0243,-0.7747) (-0.0000,-0.7746) (0.0243,-0.7747) (0.0488,-0.7750) (0.0733,-0.7755) (0.0980,-0.7761) (0.1231,-0.7770) (0.1484,-0.7780) (0.1742,-0.7792) (0.2004,-0.7806) (0.2272,-0.7821) (0.2547,-0.7837) (0.2828,-0.7855) (0.3118,-0.7874) (0.3416,-0.7893) (0.3724,-0.7913) (0.4042,-0.7933) (0.4371,-0.7951) (0.4712,-0.7968) (0.5066,-0.7982) (0.5432,-0.7993) (0.5812,-0.7999) (0.6205,-0.7999) (0.6611,-0.7991) (0.7030,-0.7974) (0.7461,-0.7945) (0.7903,-0.7903) (0.8354,-0.7845) (0.8814,-0.7770) (0.9279,-0.7677) (0.9749,-0.7562) (1.0220,-0.7425) (1.0689,-0.7264) (1.1155,-0.7079) (1.1613,-0.6868) (1.2063,-0.6632) (1.2500,-0.6369) (1.2922,-0.6081) (1.3327,-0.5767) (1.3713,-0.5429) (1.4076,-0.5068) (1.4415,-0.4684) (1.4727,-0.4279) (1.5012,-0.3854) (1.5267,-0.3412) (1.5490,-0.2955) (1.5682,-0.2484) (1.5840,-0.2001) (1.5964,-0.1509) (1.6053,-0.1010) (1.6107,-0.0506) (1.6125,-0.0000)};
\draw[structurecolor] plot coordinates {(1.8708,0.0000) (1.8692,0.0587) (1.8642,0.1173) (1.8559,0.1754) (1.8444,0.2330) (1.8297,0.2898) (1.8119,0.3456) (1.7911,0.4004) (1.7674,0.4538) (1.7408,0.5058) (1.7116,0.5561) (1.6799,0.6048) (1.6458,0.6516) (1.6096,0.6965) (1.5713,0.7394) (1.5312,0.7802) (1.4894,0.8188) (1.4462,0.8553) (1.4018,0.8896) (1.3562,0.9217) (1.3098,0.9516) (1.2626,0.9794) (1.2150,1.0051) (1.1669,1.0288) (1.1187,1.0505) (1.0703,1.0703) (1.0221,1.0884) (0.9740,1.1048) (0.9263,1.1196) (0.8788,1.1330) (0.8319,1.1450) (0.7854,1.1557) (0.7395,1.1653) (0.6942,1.1738) (0.6494,1.1813) (0.6053,1.1880) (0.5618,1.1938) (0.5189,1.1990) (0.4765,1.2035) (0.4347,1.2074) (0.3934,1.2108) (0.3526,1.2138) (0.3123,1.2163) (0.2723,1.2184) (0.2328,1.2202) (0.1935,1.2216) (0.1545,1.2228) (0.1157,1.2236) (0.0770,1.2243) (0.0385,1.2246) (0.0000,1.2247) (-0.0385,1.2246) (-0.0770,1.2243) (-0.1157,1.2236) (-0.1545,1.2228) (-0.1935,1.2216) (-0.2328,1.2202) (-0.2723,1.2184) (-0.3123,1.2163) (-0.3526,1.2138) (-0.3934,1.2108) (-0.4347,1.2074) (-0.4765,1.2035) (-0.5189,1.1990) (-0.5618,1.1938) (-0.6053,1.1880) (-0.6494,1.1813) (-0.6942,1.1738) (-0.7395,1.1653) (-0.7854,1.1557) (-0.8319,1.1450) (-0.8788,1.1330) (-0.9263,1.1196) (-0.9740,1.1048) (-1.0221,1.0884) (-1.0703,1.0703) (-1.1187,1.0505) (-1.1669,1.0288) (-1.2150,1.0051) (-1.2626,0.9794) (-1.3098,0.9516) (-1.3562,0.9217) (-1.4018,0.8896) (-1.4462,0.8553) (-1.4894,0.8188) (-1.5312,0.7802) (-1.5713,0.7394) (-1.6096,0.6965) (-1.6458,0.6516) (-1.6799,0.6048) (-1.7116,0.5561) (-1.7408,0.5058) (-1.7674,0.4538) (-1.7911,0.4004) (-1.8119,0.3456) (-1.8297,0.2898) (-1.8444,0.2330) (-1.8559,0.1754) (-1.8642,0.1173) (-1.8692,0.0587) (-1.8708,0.0000) (-1.8692,-0.0587) (-1.8642,-0.1173) (-1.8559,-0.1754) (-1.8444,-0.2330) (-1.8297,-0.2898) (-1.8119,-0.3456) (-1.7911,-0.4004) (-1.7674,-0.4538) (-1.7408,-0.5058) (-1.7116,-0.5561) (-1.6799,-0.6048) (-1.6458,-0.6516) (-1.6096,-0.6965) (-1.5713,-0.7394) (-1.5312,-0.7802) (-1.4894,-0.8188) (-1.4462,-0.8553) (-1.4018,-0.8896) (-1.3562,-0.9217) (-1.3098,-0.9516) (-1.2626,-0.9794) (-1.2150,-1.0051) (-1.1669,-1.0288) (-1.1187,-1.0505) (-1.0703,-1.0703) (-1.0221,-1.0884) (-0.9740,-1.1048) (-0.9263,-1.1196) (-0.8788,-1.1330) (-0.8319,-1.1450) (-0.7854,-1.1557) (-0.7395,-1.1653) (-0.6942,-1.1738) (-0.6494,-1.1813) (-0.6053,-1.1880) (-0.5618,-1.1938) (-0.5189,-1.1990) (-0.4765,-1.2035) (-0.4347,-1.2074) (-0.3934,-1.2108) (-0.3526,-1.2138) (-0.3123,-1.2163) (-0.2723,-1.2184) (-0.2328,-1.2202) (-0.1935,-1.2216) (-0.1545,-1.2228) (-0.1157,-1.2236) (-0.0770,-1.2243) (-0.0385,-1.2246) (-0.0000,-1.2247) (0.0385,-1.2246) (0.0770,-1.2243) (0.1157,-1.2236) (0.1545,-1.2228) (0.1935,-1.2216) (0.2328,-1.2202) (0.2723,-1.2184) (0.3123,-1.2163) (0.3526,-1.2138) (0.3934,-1.2108) (0.4347,-1.2074) (0.4765,-1.2035) (0.5189,-1.1990) (0.5618,-1.1938) (0.6053,-1.1880) (0.6494,-1.1813) (0.6942,-1.1738) (0.7395,-1.1653) (0.7854,-1.1557) (0.8319,-1.1450) (0.8788,-1.1330) (0.9263,-1.1196) (0.9740,-1.1048) (1.0221,-1.0884) (1.0703,-1.0703) (1.1187,-1.0505) (1.1669,-1.0288) (1.2150,-1.0051) (1.2626,-0.9794) (1.3098,-0.9516) (1.3562,-0.9217) (1.4018,-0.8896) (1.4462,-0.8553) (1.4894,-0.8188) (1.5312,-0.7802) (1.5713,-0.7394) (1.6096,-0.6965) (1.6458,-0.6516) (1.6799,-0.6048) (1.7116,-0.5561) (1.7408,-0.5058) (1.7674,-0.4538) (1.7911,-0.4004) (1.8119,-0.3456) (1.8297,-0.2898) (1.8444,-0.2330) (1.8559,-0.1754) (1.8642,-0.1173) (1.8692,-0.0587) (1.8708,-0.0000)};
\draw[structurecolor] plot coordinates {(2.1448,0.0000) (2.1431,0.0674) (2.1382,0.1345) (2.1300,0.2013) (2.1186,0.2676) (2.1040,0.3332) (2.0863,0.3980) (2.0655,0.4617) (2.0418,0.5243) (2.0153,0.5855) (1.9860,0.6453) (1.9541,0.7035) (1.9198,0.7601) (1.8831,0.8149) (1.8442,0.8678) (1.8032,0.9188) (1.7604,0.9678) (1.7158,1.0147) (1.6696,1.0596) (1.6220,1.1023) (1.5731,1.1430) (1.5232,1.1815) (1.4722,1.2179) (1.4205,1.2523) (1.3680,1.2846) (1.3150,1.3150) (1.2615,1.3434) (1.2077,1.3699) (1.1537,1.3946) (1.0996,1.4176) (1.0454,1.4388) (0.9912,1.4585) (0.9371,1.4766) (0.8831,1.4933) (0.8294,1.5086) (0.7758,1.5226) (0.7225,1.5353) (0.6694,1.5469) (0.6166,1.5573) (0.5640,1.5667) (0.5118,1.5750) (0.4597,1.5825) (0.4080,1.5890) (0.3564,1.5946) (0.3051,1.5995) (0.2540,1.6035) (0.2030,1.6067) (0.1521,1.6093) (0.1014,1.6110) (0.0507,1.6121) (0.0000,1.6125) (-0.0507,1.6121) (-0.1014,1.6110) (-0.1521,1.6093) (-0.2030,1.6067) (-0.2540,1.6035) (-0.3051,1.5995) (-0.3564,1.5946) (-0.4080,1.5890) (-0.4597,1.5825) (-0.5118,1.5750) (-0.5640,1.5667) (-0.6166,1.5573) (-0.6694,1.5469) (-0.7225,1.5353) (-0.7758,1.5226) (-0.8294,1.5086) (-0.8831,1.4933) (-0.9371,1.4766) (-0.9912,1.4585) (-1.0454,1.4388) (-1.0996,1.4176) (-1.1537,1.3946) (-1.2077,1.3699) (-1.2615,1.3434) (-1.3150,1.3150) (-1.3680,1.2846) (-1.4205,1.2523) (-1.4722,1.2179) (-1.5232,1.1815) (-1.5731,1.1430) (-1.6220,1.1023) (-1.6696,1.0596) (-1.7158,1.0147) (-1.7604,0.9678) (-1.8032,0.9188) (-1.8442,0.8678) (-1.8831,0.8149) (-1.9198,0.7601) (-1.9541,0.7035) (-1.9860,0.6453) (-2.0153,0.5855) (-2.0418,0.5243) (-2.0655,0.4617) (-2.0863,0.3980) (-2.1040,0.3332) (-2.1186,0.2676) (-2.1300,0.2013) (-2.1382,0.1345) (-2.1431,0.0674) (-2.1448,0.0000) (-2.1431,-0.0674) (-2.1382,-0.1345) (-2.1300,-0.2013) (-2.1186,-0.2676) (-2.1040,-0.3332) (-2.0863,-0.3980) (-2.0655,-0.4617) (-2.0418,-0.5243) (-2.0153,-0.5855) (-1.9860,-0.6453) (-1.9541,-0.7035) (-1.9198,-0.7601) (-1.8831,-0.8149) (-1.8442,-0.8678) (-1.8032,-0.9188) (-1.7604,-0.9678) (-1.7158,-1.0147) (-1.6696,-1.0596) (-1.6220,-1.1023) (-1.5731,-1.1430) (-1.5232,-1.1815) (-1.4722,-1.2179) (-1.4205,-1.2523) (-1.3680,-1.2846) (-1.3150,-1.3150) (-1.2615,-1.3434) (-1.2077,-1.3699) (-1.1537,-1.3946) (-1.0996,-1.4176) (-1.0454,-1.4388) (-0.9912,-1.4585) (-0.9371,-1.4766) (-0.8831,-1.4933) (-0.8294,-1.5086) (-0.7758,-1.5226) (-0.7225,-1.5353) (-0.6694,-1.5469) (-0.6166,-1.5573) (-0.5640,-1.5667) (-0.5118,-1.5750) (-0.4597,-1.5825) (-0.4080,-1.5890) (-0.3564,-1.5946) (-0.3051,-1.5995) (-0.2540,-1.6035) (-0.2030,-1.6067) (-0.1521,-1.6093) (-0.1014,-1.6110) (-0.0507,-1.6121) (-0.0000,-1.6125) (0.0507,-1.6121) (0.1014,-1.6110) (0.1521,-1.6093) (0.2030,-1.6067) (0.2540,-1.6035) (0.3051,-1.5995) (0.3564,-1.5946) (0.4080,-1.5890) (0.4597,-1.5825) (0.5118,-1.5750) (0.5640,-1.5667) (0.6166,-1.5573) (0.6694,-1.5469) (0.7225,-1.5353) (0.7758,-1.5226) (0.8294,-1.5086) (0.8831,-1.4933) (0.9371,-1.4766) (0.9912,-1.4585) (1.0454,-1.4388) (1.0996,-1.4176) (1.1537,-1.3946) (1.2077,-1.3699) (1.2615,-1.3434) (1.3150,-1.3150) (1.3680,-1.2846) (1.4205,-1.2523) (1.4722,-1.2179) (1.5232,-1.1815) (1.5731,-1.1430) (1.6220,-1.1023) (1.6696,-1.0596) (1.7158,-1.0147) (1.7604,-0.9678) (1.8032,-0.9188) (1.8442,-0.8678) (1.8831,-0.8149) (1.9198,-0.7601) (1.9541,-0.7035) (1.9860,-0.6453) (2.0153,-0.5855) (2.0418,-0.5243) (2.0655,-0.4617) (2.0863,-0.3980) (2.1040,-0.3332) (2.1186,-0.2676) (2.1300,-0.2013) (2.1382,-0.1345) (2.1431,-0.0674) (2.1448,-0.0000)};
\node at (-1,0) {$\otimes$};\node at (1,0) {$\otimes$};
\draw[->,thick] (-1.4,.65)--(-1.1,.75);\draw[->,thick] (1.1,.75)--(1.4,.65);\draw[->,thick] (-.2,2.14)--(.2,2.14);
\end{tikzpicture}\end{center}

两根$\otimes$表示电流沿$-x$入图, 图上的方向近线为顺时针. 场线是$[(y-a)^2+z^2][(y+a)^2+z^2]$的常值线: 对其对数函数求梯度可核验磁场处处相切. 圆形只在极近单线或极远两线的近似中出现; 中间形状由叠加而变. 图中的曲线由这条解析等值条件采样, 不以图作为极值求导的证明.

(b)$I_{\rm enc}(s)=2\pi\alpha\int_0^su^2\dd u=2\pi\alpha s^3/3$, 总$I=2\pi\alpha R^3/3$. 故$B_\phi=\mu_0\alpha s^2/3$($s<R$), 外部$\mu_0\alpha R^3/(3s)$. 场在边界连续, 内部是二次而非均匀电流的线性增长.

(c)小带电流$\dd I=I\dd u/w$, 轴距$s+u$, 所以
\[B_z=\frac{\mu_0I}{2\pi w}\int_0^w\frac{\dd u}{s+u}=\frac{\mu_0I}{2\pi w}\ln\left(1+\frac ws\right).\]
此处点在$+y$一侧, 场向$+z$; 另一侧反向. $w\ll s$时回到细线公式. 贴近无限薄理想带边缘出现对数发散, 实际厚度会改变该极限.
\end{solution}

一个低速运动点电荷在近似匀速、非相对论范围的场为
\[\bb\simeq\frac{\mu_0q}{4\pi}\frac{\vct v\times\vct R}{R^3}=\frac1{c^2}\vct v\times\ee.\]
这和许多电荷组成的稳态毕奥--萨伐尔式相容. 快速加速电荷有延迟和辐射场, 单凭瞬时位置速度的式子不完整. 一枚电荷在环上转动的瞬时场也非处处轴对称; 均匀连续带电环稳定转动或完整时间平均才得到环电流模型.

半径$R$、均匀面密度$\sigma$的圆盘以角速度$\Omega$绕轴转动, 一小环电荷$\dd Q=2\pi\sigma s\dd s$, 周期$2\pi/\Omega$, 环电流$\dd I=\sigma\Omega s\dd s$. 因此中心场与磁矩分别为
\[B_z(0)=\int_0^R\frac{\mu_0\dd I}{2s}=\frac{\mu_0\sigma\Omega R}2,\quad
m_z=\int_0^R\pi s^2\dd I=\frac{\pi\sigma\Omega R^4}4=\frac{Q\Omega R^2}4.\]
这是原9.13.5中的经典模型, 要求维持给定电荷分布且边缘速度$\Omega R\ll c$. 单个轨道电荷的磁矩是$q\vct L/(2m)$; 电子由于$q=-e$而与轨道角动量反向. 此结论不包括自旋磁矩.

\originalsection{有限螺线管与离轴圆环}\label{tb:offaxis-loop}
原例9.2还求固定轴向小磁矩$m_z$所受力. 由圆环轴场求导,
\[F_z=m_z\frac{\dd B_z}{\dd z}=-\frac{3\mu_0m_zIR^2z}{2(R^2+z^2)^{5/2}}.\]
$z>0,m_zI>0$时力朝环心, 同向吸引; 固定反向磁矩则排斥. 若磁矩可以自由转动, 还要计力矩而不能保持反向取向假定.

原例9.4的均匀厚电流板沿$+x$导电, $J_x=J_0$占$|z|<b/2$. 平移对称和反射使场只沿$y$, 叠加薄面或安培矩形给
\[B_y(z)=\begin{cases}+\mu_0J_0b/2,&z<-b/2,\\-\mu_0J_0z,&|z|<b/2,\\-\mu_0J_0b/2,&z>b/2.\end{cases}\]
$K=J_0b$单位A/m, 固定$K$令$b\to0$才还原薄电流面. 原例9.5理想密绕环形管则孔内$r<a$和环外$r>b$近似场零, 仅$a<r<b$环体中有$\mu_0NI/(2\pi r)$; 在实际有限截面还要限制高度位于环体中. 内孔环路不链接匝电流, 外环路穿过每匝去回两次而净电流零, 加上环形理想对称性才可求零场.

有限螺线管中心在原点, 半径$R$, 长$\ell$, 绕组密度$n$. 把每段看作$n\dd z'$个圆环, 轴上
\begin{align*}
B_z(z)&=\frac{\mu_0nI R^2}2\int_{-\ell/2}^{\ell/2}\frac{\dd z'}{[R^2+(z-z')^2]^{3/2}}\\
&=\frac{\mu_0nI}2\left[\frac{z+\ell/2}{\sqrt{R^2+(z+\ell/2)^2}}-\frac{z-\ell/2}{\sqrt{R^2+(z-\ell/2)^2}}\right].
\end{align*}
长管中心趋近$\mu_0nI$, 端点趋近其一半. 轴外有限管外场一般不零, 因而不能用无限管内外跳变的图代替实际端部场.

原9.8附录计算离轴圆环. 令源点为$(R\cos\phi',R\sin\phi',0)$, 观测点为$(s,0,z)$, 定义$D^2=R^2+s^2+z^2-2Rs\cos\phi'$. 叉积直接给出
\[B_s=\frac{\mu_0I}{4\pi}\int_0^{2\pi}\frac{Rz\cos\phi'}{D^3}\dd\phi',\quad
B_z=\frac{\mu_0I}{4\pi}\int_0^{2\pi}\frac{R^2-Rs\cos\phi'}{D^3}\dd\phi'.\]
$B_\phi=0$是对称性的结果. $s=0$时第一式为零、第二式还原轴场. 这些一维积分本身已给精确稳态模型, 一般离轴值需数值积分或椭圆积分, 不应硬套轴上闭式.

在$r=\sqrt{s^2+z^2}\gg R$处展开$D^{-3}$, 保留成环后第一个非零阶. 常数项使$B_s$的余弦积分为零; 一阶项贡献$3R^2sz/r^5$乘$\pi$. 于是以$m=I\pi R^2$记,
\[B_s\simeq\frac{\mu_0m}{4\pi r^3}\frac{3sz}{r^2},\quad
B_z\simeq\frac{\mu_0m}{4\pi r^3}\left(\frac{3z^2}{r^2}-1\right),\]
即$\bb=\mu_0[3(\vct m\cdot\uv r)\uv r-\vct m]/(4\pi r^3)$. 在原点附近这不是完整圆环场, 更不能拿远场发散作为环中心真实奇点.

原9.9的两线圈半径$R$, 中心在$z=\pm d/2$, 电流同向. 场是两个轴场相加, 中心一阶导数为零, 而
\[B_z''(0)=\frac{3\mu_0IR^2(d^2-R^2)}{(R^2+d^2/4)^{7/2}}.\]
令$d=R$消去二阶变化, 即Helmholtz条件, 中心场$B_0=(4/5)^{3/2}\mu_0I/R$. 同向线圈的中心平坦由导数检验, 而不是因为画出的场线平行就宣称全空间匀场. 反向线圈场为两项之差, 中心场零而轴向梯度非零; 磁体在其中的振动或悬浮需结合第\ref{ch:magneticforce}章机械势和稳定性条件.

偏心柱孔是安培对称性与叠加的迁移. 无限实圆柱均匀$\jj$沿轴, 内场可写$\bb=\mu_0\jj\times\rr_\perp/2$. 挖去完全在内部的平行柱孔, 孔心横向位移$\vct a$, 等于叠加密度$-\jj$的小柱. 孔内
\[\bb_{\rm hole}=\tfrac{\mu_0}2\jj\times\rr_\perp-\tfrac{\mu_0}2\jj\times(\rr_\perp-\vct a)=\tfrac{\mu_0}2\jj\times\vct a.\]
原9.13.3因而有均匀孔场, 但整体失去轴对称, 单个偏心安培圆不能直接从包围电流求点场. 它和第\ref{tb:charged-sphere}节绝缘球空腔的向量减法相连.

真实螺旋绕线还向轴向输送电流. 密绕圆柱半径$R$, 匝密度$n$, 单根连续螺旋电流$I$的面平均分量为$K_\phi=nI$, $K_z=I/(2\pi R)$. 第一分量产生近似长管轴场, 第二分量使包围整个绕组的安培圆仍链接轴向电流$I$, 环外另有$B_\phi\simeq\mu_0I/(2\pi s)$. 要求导线真实闭合时, 返回引线还会贡献场; 原9.13.4不能用理想“纯环向薄绕组外场零”消除这一额外分量.
